\scanpage{104}
Pour \unsure{que} \deleted{strate}\inserted{comp. conn.} $T_\alpha$ de $X'_T$, considérons
\[
\widehat\varphi_\alpha:\widehat\pi_1(X'_{\overline K},T_\alpha)\longrightarrow\widehat\pi_1(X,T)
\]
(c'est un hom injectif si $X'$ rev. étale de $X$, et alors les \deleted{\unsure{différentes im.}} images des $\varphi_\alpha$ sont les conjugués de l'image de l'un d'eux). Donc \inserted{si \deleted{\unsure{on peut}} la condition sur $x'_\alpha,y'_\alpha$ est vraie pour tt $\alpha$,} la formule 54 implique ``la'' condition intrigante
\[
(55)\qquad g^T_{y,x}(u)\in\bigcap_{\alpha\in\pi_0(X'_T)}\operatorname{Im}\widehat\varphi_\alpha
\qquad\forall u\in\Gamma=\Gamma_{\overline K/K}
\]
\margin{Si $X'$ est un rev. étale de $X$, l'ens \inserted{$E=\pi_0(X'_T)$} discret \deleted{est} d'une ``fibre'' qui décrit $X'$ en termes de $\pi=\pi_1(X,T)$, via le ``foncteur-fibre'' $\mathcal E$ défini par $T\to X$ ($T$ 1-connexe) -- une action de $\pi$ sur $E$, et le 2° mb de (55), est le sous-groupe \inserted{dist.} de $\pi$ noyau de $\pi\longrightarrow\mathfrak S_E$}

C'est là une condition arithmétique, nécessaire pour que le revêt \inserted{étale} fini topologique $X'$ de $X$ (p. ex se borner à ce cas) puisse \deleted{se représenter} être muni d'une \deleted{$K$-} structure de $K$-schéma, pour laquelle $f:X'\to X$ soit rat$ /K$, et telle que \deleted{$\unsure{x',y'|x,y}$}\inserted{$\forall\alpha\in\pi_0(X'_T)$,} $\forall u\in\Gamma$, $u(x'_\alpha)$ et $u(y'_\alpha)$ soient $T$-adjacents (Ce qui est une condition nettement plus faible que celle que $x',y'\in X'_K$ \deleted{)} \inserted{i.e. $x',y'$} fixes sous $\Gamma$)\ldots
\mspage{103}
